\subsection{Latent Representations Learned by CNNs}\label{sec:latent-representations-learned-by-cnns}

\subsubsection{The CNN Latent Representation}\label{sec:the-cnn-latent-representation}

As we have seen on \autopageref{sec:feature-extraction-network-and-classifier}, \hl{a CNN does not only produce a class prediction; it also learns a numerical representation of the image.} We also called this representation the \textbf{latent representation} or \textbf{data-driven feature vector} of the image, and placed it in the middle of the network, between the Feature Extraction Network (FEN, which is the convolutional part of the network) and the final MLP classifier (which is the dense part of the network).

\highspace
\begin{flushleft}
    \textcolor{Green3}{\faIcon{question-circle} \textbf{Where is the Latent Representation?}}
\end{flushleft}
Conceptually, the latent representation is located at the end of the convolutional part of the network, right before the first dense layer of the MLP classifier. In practice, it is located at the output of the last convolutional layer, after flattening:
\begin{equation*}
    \text{Image} \to \underbrace{\text{FEN}}_{\text{Conv} + \text{Pooling}} \to \boxed{\text{Latent Representation}} \to \underbrace{\text{MLP classifier}}_{\text{task-specific}} \to \text{class}
\end{equation*}
In other words, the latent representation is the \textbf{feature vector produced by the CNN before the final classifier}. Mathematically, for a generic image $I$, we can write:
\begin{equation*}
    \mathbf{z} = \mathrm{FEN}(I)
\end{equation*}
Where $\mathbf{z}$ is the latent representation of the original image $I$, and $\mathrm{FEN}(I)$ is the convolutional feature extraction network applied to the image $I$.

\highspace
\begin{flushleft}
    \textcolor{Green3}{\faIcon{not-equal} \textbf{The FEN and the classifier have different roles}}
\end{flushleft}
We reinforce the idea that we have already seen on \autopageref{sec:feature-extraction-network-and-classifier}:
\begin{equation*}
    \underbrace{\text{Image} \to \text{Convolutional layers}}_{\text{Feature Extraction Network (FEN)}} \to \mathbf{x} \to \underbrace{\text{MLP}}_{\text{Classifier}}
\end{equation*}
The FEN is the part that \textbf{extracts high-level features from pixels}, while the MLP performs feature classification.
\begin{itemize}
    \item FEN: ``\emph{What visual information is present?}''
    \item Classifier: ``\emph{Which class does this feature vector correspond to?}''
\end{itemize}
So it means the feature vector itself can be useful \textbf{even without the original classifier}.

\highspace
Thus, instead of representing an image by its raw pixels:
\begin{equation*}
    I = \left[x_1, x_2, \ldots, x_n\right]
\end{equation*}
We can represent it by its latent representation:
\begin{equation*}
    \mathbf{z} = \mathrm{FEN}(I)
\end{equation*}
Where $\mathbf{z}$ summarizes the high-level information learned by the CNN from the original image $I$.

\highspace
\begin{flushleft}
    \textcolor{Green3}{\faIcon{question-circle} \textbf{Why is this representation useful?}}
\end{flushleft}
A natural question about the latent representation is: \emph{why is it useful to have a feature vector that summarizes the image?} The answer is that \hl{images with similar semantic content tend to have similar learned representations!} Mathematically, if two images $I_1$ and $I_2$ are similar, then their latent representations $\mathbf{z}_1$ and $\mathbf{z}_2$ will also be similar:
\begin{equation}\label{eq:latent-space-distance}
    d\left(I_1 , I_2\right) = \left\| \mathrm{FEN}\left(I_1\right) - \mathrm{FEN}\left(I_2\right) \right\|_{2} = \left\| \mathbf{z}_1 - \mathbf{z}_2 \right\|_2
\end{equation}
Latent-space distance measures \textbf{how similar the features learned by the CNN are}. This is in contrast to pixel distance, which measures the similarity of raw intensity values. In other words, the distance in the latent representation of a CNN are much more meaningful than distances in the data itself.

\highspace
\begin{takeawaysbox}
    \textbf{A CNN learns not only a classifier, but also a feature representation of the input image.} The Feature Extraction Network transforms raw pixels into a \textbf{latent, data-driven feature vector}:
    \begin{equation*}
        \mathbf{z} = \mathrm{FEN}(I)
    \end{equation*}
    This vector summarizes high-level visual information and can be used independently of the final classifier. Distances between latent vectors are often more meaningful than distances between raw images because the representation has been learned to capture task-relevant visual structure.
\end{takeawaysbox}